\chapter{Evaluation and Results}
This chapter presents the evaluation methodology, experimental results, and discussion used to assess the RAG pipeline subsystem and its end-to-end performance. The comparative evaluation centers on the RAG pipeline as this is the subsystem that can be most directly assessed against the existing chatbot in terms of performance, scalability, and latency.

\section{Evaluation Pipeline Implementation}
The evaluation pipeline was designed to assess both retrieval quality and end-to-end answer quality while remaining lightweight enough to support repeated experimentation during system development. Detailed prompt templates, schemas, implementation screenshots, and extended worked examples are provided in Appendix B and Appendix C.

\subsection{Dataset Generation}
The evaluation dataset was generated automatically using an LLM-based process rather than fully manual annotation. Manual construction of question-answer benchmarks across long meeting transcripts would have been slow, costly, and difficult to scale, especially when the retrieval pipeline was still being tuned \cite{ares2023}. Recent work has shown that LLM-assisted evaluation can substantially reduce this overhead while still supporting systematic comparison \cite{ragas2023}.

To create the dataset, each transcript was first grouped into larger context windows and then passed to an LLM to synthesize grounded question-answer pairs. Each QA item retained the generated question and answer together with the source start timestamp, source end timestamp, and originating context text. The timestamp metadata was particularly important because it allowed retrieval relevance to be checked algorithmically through temporal overlap rather than manual judging. Appendix B documents the grouping example, generation prompt, API-chain setup, and ground-truth JSON structure used in this process.

Importantly, the generated dataset included a mixture of narrow fact-extraction questions and higher-level, less localized questions, such as asking for a summary of the meeting. This broader question distribution was intentional to reduce bias toward RAG-style retrieval evaluation, since a benchmark made up only of highly specific span-level questions would naturally favor systems optimized for retrieving small, precise evidence windows. Including both question types therefore provides a more balanced assessment of how the system performs across realistic transcript-query patterns.

\subsection{Evaluation Metrics}
The pipeline reports two retrieval metrics and two generation metrics so that failures can be attributed to the correct stage of the RAG system.

\textbf{Hit Rate} measures the proportion of questions for which at least one retrieved chunk overlaps the ground-truth timestamp range \cite{vanrijsbergen1979ir}. It is used as a recall-oriented indicator of whether the retriever surfaced any relevant evidence at all.

\textbf{Mean Reciprocal Rank (MRR)} measures how highly the first relevant chunk is ranked \cite{wikipedia2024mrr}. It complements Hit Rate by distinguishing between systems that eventually retrieve relevant evidence and systems that place that evidence near the top of the result list.

\textbf{F1 Score} measures token-level overlap between the generated answer and the reference answer \cite{vanrijsbergen1979ir}. It captures lexical fidelity, but it may undervalue correct paraphrases or grounded elaborations.

\textbf{LLM-Judge Accuracy} uses a separate LLM to assess whether the generated answer is semantically consistent with the reference answer \cite{liu2023geval}. This metric captures semantic correctness more directly than F1, though it remains sensitive to the design of the judge prompt and the specificity of the ground truth.

Retrieval and generation were evaluated separately for diagnostic clarity. Hit Rate and MRR are computed before answer generation and therefore isolate retrieval behavior, while F1 and LLM-judge accuracy are computed after generation and therefore reflect how well the language model used the retrieved evidence. This separation made it possible to distinguish retrieval failures from answer-generation failures during experimentation. Appendix B contains the temporal-overlap logic, F1 preprocessing details, and judge prompt used in the implementation, while Appendix C contains extended diagnostic examples.

\subsection{Speed Benchmarking Design}
Latency was evaluated using a separate speed benchmarking pipeline so that runtime comparisons between the RAG system and a Pure LLM baseline would remain fair. Both methods were tested on the same question set derived from the evaluation dataset, ensuring that quality and latency analyses were based on comparable workloads.

The benchmark used interleaved execution rather than evaluating one system fully before the other. This design reduced bias from prompt caching in modern LLM services, which could otherwise make later runs appear artificially faster \cite{openai_prompt_caching}. A short warmup phase was also applied before timed runs to reduce cold-start effects from model loading, runtime initialization, and connection setup, with warmup executions excluded from the final statistics. Appendix B provides the interleaving illustration together with the supporting latency-statistics and comparative-metrics tables.

\section{Experimental Results}
This section presents the empirical results from evaluating the RAG system across multiple transcripts, covering both quality metrics (retrieval and generation performance) and speed performance (comparison with the Pure LLM baseline).

The reported experiments were conducted using a locally hosted development LLM rather than the self-hosted model deployed on the Meadow9+ production server. As a result, the absolute latency and generation quality reported here should be interpreted as development-stage measurements rather than final production benchmarks.

\subsection{RAG Quality Metrics}
The RAG quality evaluation indicates that the proposed retrieval pipeline performs reliably across most of the evaluated question set. In general, the system was able to retrieve the correct supporting evidence for the majority of queries, and end-to-end answer quality was also reasonably strong given the difficulty of evaluating grounded natural-language answers. Taken together, these results suggest that the core retrieval design is effective and that the main remaining limitations lie less in evidence discovery and more in downstream answer generation and evaluation strictness.

Table \ref{tab:rag-quality-metrics} presents the per-transcript RAG quality metrics.

\begin{table}[H]
  \centering
  \begingroup
  \small
  \setlength{\tabcolsep}{4pt}
  \begin{tabular}{|c|c|c|c|c|c|c|}
    \hline
    \textbf{Transcript} & \textbf{Duration} & \textbf{Questions} & \textbf{Hit Rate} & \textbf{MRR} & \textbf{Accuracy} & \textbf{F1 Score} \\
    \hline
    Transcript 1 & 00:08:21 & 29 & 1.000 & 1.000 & 1.000 & 0.677 \\
    \hline
    Transcript 2 & 01:59:40 & 150 & 0.867 & 0.867 & 0.767 & 0.476 \\
    \hline
    Transcript 3 & 00:08:14 & 32 & 1.000 & 1.000 & 0.875 & 0.613 \\
    \hline
    Transcript 4 & 00:36:14 & 62 & 1.000 & 1.000 & 0.800 & 0.446 \\
    \hline
    Transcript 5 & 00:34:08 & 59 & 1.000 & 1.000 & 0.900 & 0.512 \\
    \hline
    Transcript 6 & 00:36:14 & 58 & 1.000 & 1.000 & 0.733 & 0.478 \\
    \hline
    \textbf{Aggregate} & \textbf{04:02:51} & \textbf{390} & \textbf{0.949} & \textbf{0.949} & \textbf{0.775} & \textbf{0.503} \\
    \hline
  \end{tabular}
  \endgroup
  \caption{Per-transcript RAG quality metrics aggregated results}
  \label{tab:rag-quality-metrics}
\end{table}

\subsubsection*{Aggregate Interpretation}
\textbf{Retrieval performance.} Under the six-transcript evaluation set, the system achieved a question-weighted Hit Rate of 0.949 and an MRR of 0.949. This indicates that relevant evidence was retrieved and ranked highly for most questions, although Transcript 2 remains the main source of retrieval misses.

\textbf{Generation performance.} Using the per-transcript values, the system achieved a question-weighted aggregate accuracy of 77.5\% and an average F1 score of 0.503. This remains consistent with the earlier observation that retrieval is generally strong, while end-to-end answer quality is more sensitive to generation behavior and judge strictness.

\subsubsection*{False Negatives from Semantic Judge}
One recurring evaluation issue was false negatives from the semantic judge. In some cases, the generated answer was more specific than the reference answer while still remaining grounded in the transcript, causing the judge to mark it as incorrect despite the answer being substantively valid. To reduce the impact of such cases during analysis, contradictory judgments between the RAG and Pure LLM outputs were reviewed more carefully through a secondary manual-assisted process: the transcript together with both candidate answers was provided to an LLM for closer comparison, allowing obviously grounded but more specific answers to be identified before interpretation of the final results. A fuller example and additional diagnostic cases are provided in Appendix C.

Overall, the quality evaluation supports the conclusion that the RAG subsystem is technically sound for transcript-grounded question answering. Retrieval performance is consistently high, and answer accuracy is generally good despite the inherent difficulty of judging paraphrased or more specific responses.

\subsection{Speed Performance}
This is the more operationally important part of the evaluation, because it directly measures improvement over the existing Meadow9+ chatbot. While the quality evaluation shows whether the proposed RAG pipeline can answer questions correctly, the speed evaluation determines whether the new architecture delivers a practical product advantage over the current solution in terms of responsiveness and scalability under real transcript lengths.

Table \ref{tab:speed-comparison} presents the per-transcript speed comparison.

\begin{table}[H]
  \centering
  \begingroup
  \small
  \setlength{\tabcolsep}{4pt}
  \begin{tabular}{|c|c|c|c|c|}
    \hline
    \textbf{Transcript} & \textbf{Duration} & \textbf{RAG Accuracy} & \textbf{Pure LLM Accuracy} & \textbf{Speedup Factor} \\
    \hline
    Transcript 1 & 00:08:21 & 100.00\% & 100.00\% & 0.58x \\
    \hline
    Transcript 2 & 01:59:40 & 76.67\% & 53.33\% & 8.35x \\
    \hline
    Transcript 3 & 00:08:14 & 87.50\% & 100.00\% & 0.63x \\
    \hline
    Transcript 4 & 00:36:14 & 80.00\% & 70.00\% & 2.21x \\
    \hline
    Transcript 5 & 00:34:08 & 90.00\% & 66.67\% & 2.25x \\
    \hline
    Transcript 6 & 00:36:14 & 73.33\% & 70.00\% & 2.20x \\
    \hline
  \end{tabular}
  \endgroup
  \caption{Per-transcript speed comparison showing speedup factors and end-to-end accuracy}
  \label{tab:speed-comparison}
\end{table}

\subsubsection*{Speedup vs Transcript Duration}
The speed results show a clear length-dependent pattern. On the longer transcripts, Transcript 2, Transcript 4, Transcript 5 and Transcript 6, RAG was consistently faster than the Pure LLM baseline, with speedups ranging from 2.20x to 8.35x, and also achieved equal or better accuracy. The largest gain occurred on Transcript 2, the longest transcript, where RAG was both substantially faster and more accurate, which supports the claim that retrieval becomes increasingly beneficial as transcript length grows. 

On shorter transcripts, Transcript 1 (00:08:21) and Transcript 3 (00:08:14), however, RAG was slower than the Pure LLM baseline because retrieval overhead outweighed the benefit of reducing prompt length. In these sub-10-minute cases, Pure LLM also matched or exceeded RAG in accuracy, indicating that retrieval is not always the most efficient choice for shorter inputs. 

This suggests a practical direction for future work: the system could estimate transcript size or token count and switch between a Pure LLM path and a RAG-based path based on a tuned threshold. Such a threshold could be optimized empirically so that shorter transcripts avoid unnecessary retrieval overhead while longer transcripts continue to benefit from retrieval-based context reduction.

Overall, the speed evaluation provides the clearest evidence of practical improvement over the existing Meadow9+ chatbot approach. The results show that RAG is not universally better for every transcript length, but it offers a strong advantage on the longer transcripts that are most likely to stress a pure-LLM pipeline in production. This supports the view that the proposed system is not only viable as a retrieval-augmented chatbot architecture, but also meaningfully more scalable than the current baseline when transcript size grows.
